\batchmode
\documentclass[12pt,a4paper]{article}
\RequirePackage{ifthen}



\usepackage[utf8]{inputenc}
\usepackage[T1]{fontenc}
\usepackage[brazilian]{babel}

\usepackage{geometry}
\geometry{margin=2.5cm}

\usepackage{setspace}
\onehalfspacing

\usepackage{amsmath,amssymb}
\usepackage{enumitem}
\usepackage{hyperref}
\usepackage{booktabs}
\usepackage{array}
\usepackage{xcolor}
\usepackage{tcolorbox}
\usepackage{framed}
\usepackage{graphicx}
\usepackage{longtable}

\hypersetup{
    colorlinks=true,
    linkcolor=blue,
    urlcolor=blue,
    citecolor=blue
}

\definecolor{mental}{RGB}{30,100,150}
\definecolor{social}{RGB}{200,120,30}
\definecolor{ambiental}{RGB}{30,120,50}
\definecolor{modus}{RGB}{100,50,150}
\definecolor{alerta}{RGB}{200,0,0}
\definecolor{dados}{RGB}{0,100,100}
\definecolor{wulf}{RGB}{150,50,100}

\title{WebQuest: Quanto custa o que você consome?}
\author{Assinaturas digitais: dinheiro, gente e planeta -- \\
Uma perspectiva antropológica a partir de Christoph Wulf}
\date{}





\makeatletter

\makeatletter
\count@=\the\catcode`\_ \catcode`\_=8 
\newenvironment{tex2html_wrap}{}{}%
\catcode`\<=12\catcode`\_=\count@
\newcommand{\providedcommand}[1]{\expandafter\providecommand\csname #1\endcsname}%
\newcommand{\renewedcommand}[1]{\expandafter\providecommand\csname #1\endcsname{}%
  \expandafter\renewcommand\csname #1\endcsname}%
\newcommand{\newedenvironment}[1]{\newenvironment{#1}{}{}\renewenvironment{#1}}%
\let\newedcommand\renewedcommand
\let\renewedenvironment\newedenvironment
\makeatother
\let\mathon=$
\let\mathoff=$
\ifx\AtBeginDocument\undefined \newcommand{\AtBeginDocument}[1]{}\fi
\newbox\sizebox
\setlength{\hoffset}{0pt}\setlength{\voffset}{0pt}
\addtolength{\textheight}{\footskip}\setlength{\footskip}{0pt}
\addtolength{\textheight}{\topmargin}\setlength{\topmargin}{0pt}
\addtolength{\textheight}{\headheight}\setlength{\headheight}{0pt}
\addtolength{\textheight}{\headsep}\setlength{\headsep}{0pt}
\newwrite\lthtmlwrite
\makeatletter
\let\realnormalsize=\normalsize
\global\topskip=2sp
\def\preveqno{}\let\real@float=\@float \let\realend@float=\end@float
\def\@float{\let\@savefreelist\@freelist\real@float}
\def\liih@math{\ifmmode$\else\bad@math\fi}
\def\end@float{\realend@float\global\let\@freelist\@savefreelist}
\let\real@dbflt=\@dbflt \let\end@dblfloat=\end@float
\let\@largefloatcheck=\relax
\let\if@boxedmulticols=\iftrue
\def\@dbflt{\let\@savefreelist\@freelist\real@dbflt}
\def\adjustnormalsize{\def\normalsize{\mathsurround=0pt \realnormalsize
 \parindent=0pt\abovedisplayskip=0pt\belowdisplayskip=0pt}%
 \def\phantompar{\csname par\endcsname}\normalsize}%
\def\lthtmltypeout#1{{\let\protect\string \immediate\write\lthtmlwrite{#1}}}%
\usepackage[tightpage,active]{preview}
\PreviewBorder=0.5bp
\newbox\lthtmlPageBox
\newdimen\lthtmlCropMarkHeight
\newdimen\lthtmlCropMarkDepth
\long\def\lthtmlTightVBoxA#1{\def\lthtmllabel{#1}
    \setbox\lthtmlPageBox\vbox\bgroup\catcode`\_=8 }%
\long\def\lthtmlTightVBoxZ{\egroup
    \lthtmlCropMarkHeight=\ht\lthtmlPageBox \advance \lthtmlCropMarkHeight 6pt
    \lthtmlCropMarkDepth=\dp\lthtmlPageBox
    \lthtmltypeout{^^J:\lthtmllabel:lthtmlCropMarkHeight:=\the\lthtmlCropMarkHeight}%
    \lthtmltypeout{^^J:\lthtmllabel:lthtmlCropMarkDepth:=\the\lthtmlCropMarkDepth:1ex:=\the \dimexpr 1ex}%
    \begin{preview}\copy\lthtmlPageBox\end{preview}}%
\long\def\lthtmlTightFBoxA#1{\def\lthtmllabel{#1}%
    \adjustnormalsize\setbox\lthtmlPageBox=\vbox\bgroup\hbox\bgroup %
    \let\ifinner=\iffalse \let\)\liih@math %
    \bgroup\catcode`\_=8 }%
\long\def\lthtmlTightFBoxZ{\egroup\egroup
    \@next\next\@currlist{}{\def\next{\voidb@x}}%
    \expandafter\box\next\egroup %
    \lthtmlCropMarkHeight=\ht\lthtmlPageBox \advance \lthtmlCropMarkHeight 6pt
    \lthtmlCropMarkDepth=\dp\lthtmlPageBox
    \lthtmltypeout{^^J:\lthtmllabel:lthtmlCropMarkHeight:=\the\lthtmlCropMarkHeight}%
    \lthtmltypeout{^^J:\lthtmllabel:lthtmlCropMarkDepth:=\the\lthtmlCropMarkDepth:1ex:=\the \dimexpr 1ex}%
    \begin{preview}\copy\lthtmlPageBox\end{preview}}%
    \long\def\lthtmlinlinemathA#1#2\lthtmlindisplaymathZ{\lthtmlTightVBoxA{#1}{\hbox\bgroup#2\egroup}\lthtmlTightVBoxZ}
    \def\lthtmlinlineA#1#2\lthtmlinlineZ{\lthtmlTightVBoxA{#1}{\hbox\bgroup#2\egroup}\lthtmlTightVBoxZ}
    \long\def\lthtmldisplayA#1#2\lthtmldisplayZ{\lthtmlTightVBoxA{#1}{#2}\lthtmlTightVBoxZ}
    \long\def\lthtmldisplayB#1#2\lthtmldisplayZ{\\edef\preveqno{(\theequation)}%
        \lthtmlTightVBoxA{#1}{\let\@eqnnum\relax#2}\lthtmlTightVBoxZ}
    \long\def\lthtmlfigureA#1{\let\@savefreelist\@freelist
        \lthtmlTightFBoxA{#1}}
    \long\def\lthtmlfigureZ{
        \lthtmlTightFBoxZ\global\let\@freelist\@savefreelist}
    \long\def\lthtmlpictureA#1{\let\@savefreelist\@freelist
        \lthtmlTightVBoxA{#1}}
    \long\def\lthtmlpictureZ{
        \lthtmlTightVBoxZ\global\let\@freelist\@savefreelist}
\def\lthtmlcheckvsize{\ifdim\ht\sizebox<\vsize 
  \ifdim\wd\sizebox<\hsize\expandafter\hfill\fi \expandafter\vfill
  \else\expandafter\vss\fi}%
\providecommand{\selectlanguage}[1]{}%
\makeatletter \tracingstats = 1 
\providecommand{\Alpha}{\textrm{A}}
\providecommand{\Beta}{\textrm{B}}
\providecommand{\Chi}{\textrm{X}}
\providecommand{\Epsilon}{\textrm{E}}
\providecommand{\Eta}{\textrm{H}}
\providecommand{\Iota}{\textrm{J}}
\providecommand{\Kappa}{\textrm{K}}
\providecommand{\Mu}{\textrm{M}}
\providecommand{\Nu}{\textrm{N}}
\providecommand{\Omicron}{\textrm{O}}
\providecommand{\Rho}{\textrm{R}}
\providecommand{\Tau}{\textrm{T}}
\providecommand{\Zeta}{\textrm{Z}}
\providecommand{\omicron}{\textrm{o}}


\begin{document}
\pagestyle{empty}\thispagestyle{empty}\lthtmltypeout{}%
\lthtmltypeout{latex2htmlLength hsize=\the\hsize}\lthtmltypeout{}%
\lthtmltypeout{latex2htmlLength vsize=\the\vsize}\lthtmltypeout{}%
\lthtmltypeout{latex2htmlLength hoffset=\the\hoffset}\lthtmltypeout{}%
\lthtmltypeout{latex2htmlLength voffset=\the\voffset}\lthtmltypeout{}%
\lthtmltypeout{latex2htmlLength topmargin=\the\topmargin}\lthtmltypeout{}%
\lthtmltypeout{latex2htmlLength topskip=\the\topskip}\lthtmltypeout{}%
\lthtmltypeout{latex2htmlLength headheight=\the\headheight}\lthtmltypeout{}%
\lthtmltypeout{latex2htmlLength headsep=\the\headsep}\lthtmltypeout{}%
\lthtmltypeout{latex2htmlLength parskip=\the\parskip}\lthtmltypeout{}%
\lthtmltypeout{latex2htmlLength oddsidemargin=\the\oddsidemargin}\lthtmltypeout{}%
\makeatletter
\if@twoside\lthtmltypeout{latex2htmlLength evensidemargin=\the\evensidemargin}%
\else\lthtmltypeout{latex2htmlLength evensidemargin=\the\oddsidemargin}\fi%
\lthtmltypeout{}%
\makeatother
\setcounter{page}{1}
\onecolumn

% !!! IMAGES START HERE !!!

{\newpage\clearpage
\lthtmlinlinemathA{tex2html_wrap_inline671}%
\fbox{\begin{minipage}{0.9\textwidth}
\begin{center}\textbf{\Large WebQuest: "Quanto custa o que você consome?"}
\end{center}
\par
\textbf{Título completo:} Assinaturas digitais: quanto custa, de verdade, o que você consome? -- Uma WebQuest sobre dinheiro, gente e planeta
\par
\textbf{Público-alvo:} 9º ano do Ensino Fundamental ao 2º ano do Ensino Médio
\par
\textbf{Disciplinas integradas:} Matemática, Língua Portuguesa, Ciências da Natureza, Geografia, Educação Financeira
\par
\textbf{Tempo estimado:} 4 aulas de 50 minutos
\end{minipage}}%
\lthtmlindisplaymathZ
\lthtmlcheckvsize\clearpage}

\stepcounter{section}
{\newpage\clearpage
\lthtmlinlinemathA{tex2html_wrap_inline683}%
\fbox{\begin{minipage}{0.85\textwidth}
\begin{center}\textbf{\Large Por que esta WebQuest é importante hoje?}
\end{center}
\par
Como Christoph Wulf (2022) nos ensina, \textbf{vivemos no Antropoceno, a Era dos Humanos}. É uma era marcada pela influência destrutiva dos seres humanos no planeta. Estamos subjugando a natureza, explorando-a imprudentemente e destruindo os fundamentos de nossas próprias vidas.
\par
\begin{quote}
\textit{"A educação é um componente constitutivo do ser humano; e o conhecimento do ser humano é um elemento constitutivo da educação. Os dois estão inseparavelmente entrelaçados."} (Wulf, 2022, p. 16)
\end{quote}
\par
Nesta WebQuest, você não vai apenas calcular quanto gasta com streaming. Vai \textbf{compreender como a cultura digital está transformando sua vida}, sua relação com o dinheiro, com os outros e com o planeta. Vai pensar sobre \textbf{o que significa ser humano na Era Digital.}
\end{minipage}}%
\lthtmlindisplaymathZ
\lthtmlcheckvsize\clearpage}

{\newpage\clearpage
\lthtmlinlinemathA{tex2html_wrap_inline687}%
\fbox{\begin{minipage}{0.85\textwidth}
\begin{center}\textbf{\Large A pergunta central desta WebQuest é:}
\end{center}
\textit{O que estamos fazendo com nosso tempo, dinheiro, atenção e planeta ao assinar tantos serviços?}
\end{minipage}}%
\lthtmlindisplaymathZ
\lthtmlcheckvsize\clearpage}

{\newpage\clearpage
\lthtmlinlinemathA{tex2html_wrap_inline693}%
\fbox{\begin{minipage}{0.85\textwidth}
\begin{center}\textbf{\Large E, como diria Wulf:}
\end{center}
\textit{Como a digitalização e a inteligência artificial estão modificando os fundamentos de nossas condições mundiais sociais e culturais e determinando nosso entendimento de mundo e o nosso eu?} (Wulf, 2022, p. 7)
\end{minipage}}%
\lthtmlindisplaymathZ
\lthtmlcheckvsize\clearpage}

{\newpage\clearpage
\lthtmlinlinemathA{tex2html_wrap_inline709}%
\fbox{\begin{minipage}{0.95\textwidth}
\begin{center}\textcolor{alerta}{\textbf{\Large O CONSUMO ENERGÉTICO DOS DATA CENTERS}}
\end{center}
\par
De acordo com o relatório \textbf{Electricity 2024 -- Analysis and Forecast to 2026} da Agência Internacional de Energia (IEA):
\par
\begin{tabular}{|p{4cm}|p{6cm}|}
\hline
\textbf{Indicador} & \textbf{Valor} \\
\hline
Consumo global de data centers (2022) & 460 TWh (2\% da demanda elétrica global) \\
\hline
Consumo previsto (2026) & \textcolor{alerta}{MAIS DE 1.000 TWh} \\
\hline
\textbf{Crescimento} & \textbf{MAIS QUE DOBRA em apenas 4 anos} \\
\hline
Equivalência & Todo o consumo elétrico do \textbf{JAPÃO} (3ª maior economia do mundo) \\
\hline
\end{tabular}
\par
\vspace{0.5cm}
\par
\textbf{Cenários da IEA para 2026:}
\begin{itemize}
    \item Cenário conservador: 650 TWh (Energia equivalente à Suécia)
    \item Cenário base: 1.000+ TWh (Energia equivalente ao Japão)
    \item Cenário alto: 1.050 TWh (Energia equivalente à Alemanha)
\end{itemize}
\par
\vspace{0.3cm}
\textit{Fonte: Agência Internacional de Energia (IEA) - Electricity 2024}
\end{minipage}}%
\lthtmlindisplaymathZ
\lthtmlcheckvsize\clearpage}

{\newpage\clearpage
\lthtmlinlinemathA{tex2html_wrap_inline715}%
\fbox{\begin{minipage}{0.95\textwidth}
\begin{center}\textcolor{alerta}{\textbf{\Large O MOTOR DA EXPLOSÃO: INTELIGÊNCIA ARTIFICIAL}}
\end{center}
\par
\begin{tabular}{|p{4cm}|p{6cm}|}
\hline
\textbf{Comparação} & \textbf{Consumo de energia} \\
\hline
Busca tradicional no Google & 0,3 watt-hora \\
\hline
Consulta com IA generativa (ChatGPT, Gemini, Claude) & 3 watt-hora \\
\par
\hline
\end{tabular}
\par
\vspace{0.5cm}
\par
\textbf{Fatores que empurram o consumo para cima:}
\begin{enumerate}
    \item Aumento do processamento por IA generativa
    \item Necessidade crescente de resfriamento dos servidores
    \item Nova onda de mineração de criptomoedas (160 TWh previstos para 2026)
\end{enumerate}
\par
\end{minipage}}%
\lthtmlindisplaymathZ
\lthtmlcheckvsize\clearpage}

{\newpage\clearpage
\lthtmlinlinemathA{tex2html_wrap_inline721}%
\fbox{\begin{minipage}{0.95\textwidth}
\begin{center}\textcolor{dados}{\textbf{\Large O IMPACTO POR PAÍS}}
\end{center}
\par
\textbf{ESTADOS UNIDOS:}
\begin{itemize}
    \item Concentra 33\% de todos os data centers do planeta (mais de 8.000 centros)
    \item Consumo: 200 TWh (2022) → 260 TWh (2026)
    \item Até 2030: data centers serão responsáveis por \textbf{50\% de todo o crescimento elétrico americano}
\end{itemize}
\par
\textbf{IRLANDA: o caso extremo}
\begin{itemize}
    \item 2022: data centers consumiam 17\% da eletricidade do país
    \item 2026: projeção de \textbf{32\%} -- um terço de toda a energia elétrica da Irlanda
    \item 82 data centers operacionais + 14 em construção + 40 aprovados
\end{itemize}
\par
\textbf{BRASIL:}
\begin{itemize}
    \item Amazon AWS, Microsoft Azure e Google Cloud expandindo em SP (Vinhedo, Hortolândia, Tamboré)
    \item Consumo de data centers: 304 MW (2026) → 3.457 MW (2030) = \textbf{CRESCIMENTO DE 11 VEZES}
    \item Investimento previsto: R\$ 60 bilhões até 2030 (Thymos Energia)
    \item Paradoxo: Brasil tem energia solar sobrando ao meio-dia, mas sem armazenamento
    \item Como voçe aluno armazenaria essa energia solar que sobra no Brasil ao meio-dia
\end{itemize}
\end{minipage}}%
\lthtmlindisplaymathZ
\lthtmlcheckvsize\clearpage}

{\newpage\clearpage
\lthtmlinlinemathA{tex2html_wrap_inline727}%
\fbox{\begin{minipage}{0.95\textwidth}
\begin{center}\textcolor{alerta}{\textbf{\Large A COROAÇÃO DA IA: FOME DE ENERGIA}}
\end{center}
\par
O grande catalisador da virada é a IA generativa. Antes de 2023, a previsão era que os data centers crescessem gradualmente. Mas o lançamento do ChatGPT no fim de 2022 mudou totalmente o cenário.
\par
\textbf{Projeções futuras:}
\begin{itemize}
    \item Em janeiro de 2025, a IEA atualizou seus modelos: consumo de até \textbf{1.637 TWh em 2030} no cenário mais otimista
    \item A chamada \textbf{"IA agentic"} (agentes autônomos) pode triplicar o consumo de eletricidade dos data centers existentes
    \item Empresas como Microsoft, Amazon, Google e Meta buscam fontes próprias de energia: geotermia, nuclear pequeno modular e até reabertura de usinas a carvão
\end{itemize}
\par
\textit{"A pergunta que fica é: se os data centers de IA vão consumir tanta energia quanto o Japão em 2026 e potencialmente o triplo até 2030, quem vai pagar a conta?"}
\begin{itemize}
    \item \textbf{Pergunta}: O que é energia geotermica?
\end{itemize}
\par
\end{minipage}}%
\lthtmlindisplaymathZ
\lthtmlcheckvsize\clearpage}

{\newpage\clearpage
\lthtmlinlinemathA{tex2html_wrap_inline757}%
\fbox{\begin{minipage}{0.9\textwidth}
\begin{center}\textcolor{obsolescencia}{\textbf{⚙️ O CICLO DA OBSOLESCÊNCIA PROGRAMADA}}
\end{center}
\par
\begin{tabular}{|p{2.5cm}|p{6cm}|}
\hline
\textbf{Fase} & \textbf{O que significa} \\
\hline
Duração & Produtos são feitos para durar pouco (celular lento, lâmpada que queima rápido) \\
\hline
Consumo & Você é forçado a comprar sempre, mesmo sem querer \\
\hline
Descartar & O produto velho vira lixo \\
\hline
Reciclagem & (Se acontece) consome energia e água -- mas só 20\% do lixo eletrônico é reciclado \\
\hline
Matéria-prima natural & Minérios, petróleo, madeira, água são extraídos da natureza \\
\hline
Água & Usada na produção e contamina rios \\
\hline
Energia & Usada na produção, transporte, armazenamento \\
\hline
Dióxido de carbono & Emitido na produção e transporte -- aquece o planeta \\
\hline
Exploração da matéria-prima & Mineração predatória, trabalho infantil, desmatamento (ex: coltan no Congo) \\
\hline
Descarte final & Aterro sanitário, lixão, incineração, exportação para países pobres \\
\hline
\end{tabular}
\par
\textit{"O streaming não é imaterial. Ele depende desse ciclo inteiro."} \\
\textbf{Obsolescência programada} (ou planejada) é a estratégia industrial que encurta intencionalmente a vida útil de um produto.
\par
\end{minipage}}%
\lthtmlindisplaymathZ
\lthtmlcheckvsize\clearpage}

{\newpage\clearpage
\lthtmlinlinemathA{tex2html_wrap_inline980}%
$ \rightarrow$%
\lthtmlindisplaymathZ
\lthtmlcheckvsize\clearpage}

\stepcounter{section}
\stepcounter{section}
{\newpage\clearpage
\lthtmlinlinemathA{tex2html_wrap_inline841}%
\fbox{\begin{minipage}{0.9\textwidth}
\textbf{Como descobrir:}
\begin{enumerate}
    \item Pegue a \textbf{conta de luz} da sua casa (pode ser a última que chegou).
    \item Procure o campo \textbf{"Consumo (kWh)"} ou \textbf{"Energia Elétrica (kWh)"}.
    \item Anote o valor: \_\_\_\_\_\_\_\_\_\_ kWh (consumo do mês)
\end{enumerate}
\par
\textbf{Não tem a conta em mãos?} Pergunte para um responsável:
\begin{itemize}
    \item Casa pequena (1-2 pessoas): ~150 kWh/mês
    \item Casa média (3-4 pessoas): ~250 kWh/mês
    \item Casa grande (5+ pessoas): ~400 kWh/mês
\end{itemize}
\end{minipage}}%
\lthtmlindisplaymathZ
\lthtmlcheckvsize\clearpage}

{\newpage\clearpage
\lthtmlinlinemathA{tex2html_wrap_inline855}%
\fbox{\begin{minipage}{0.9\textwidth}
\begin{center}\textbf{\Large QUANTAS CASAS CABEM EM 1.000 TWh?}
\end{center}
\par
\textbf{Fórmula:}
\begin{quote}
\textit{Número de casas = Consumo dos data centers ÷ Consumo da sua casa}
\end{quote}
\par
\textbf{Exemplo:}
\begin{itemize}
    \item Consumo data centers (2026): 1.000.000.000.000 kWh
    \item Consumo da sua casa: 250 kWh/mês = 3.000 kWh/ano
    \item \textbf{Número de casas = 1.000.000.000.000 ÷ 3.000 = 333.333.333 casas}
\end{itemize}
\par
\textbf{Agora calcule com seus números:}
\begin{itemize}
    \item Consumo data centers: \_\_\_\_\_\_\_\_\_\_ kWh
    \item Consumo da sua casa (ano): \_\_\_\_\_\_\_\_\_\_ kWh
    \item \textbf{Número de casas equivalentes: \_\_\_\_\_\_\_\_\_\_}
\end{itemize}
\par
\textit{Isso significa que os data centers consomem tanta energia quanto \textbf{333 milhões de casas} como a sua!}
\end{minipage}}%
\lthtmlindisplaymathZ
\lthtmlcheckvsize\clearpage}

{\newpage\clearpage
\lthtmlinlinemathA{tex2html_wrap_inline881}%
\fbox{\begin{minipage}{0.9\textwidth}
\begin{center}\textcolor{alerta}{\textbf{1.000 TWh = 1.000.000.000.000 kWh}}
\end{center}
\end{minipage}}%
\lthtmlindisplaymathZ
\lthtmlcheckvsize\clearpage}

{\newpage\clearpage
\lthtmlinlinemathA{tex2html_wrap_inline901}%
\fbox{\begin{minipage}{0.9\textwidth}
\begin{center}\textbf{\Large QUANTO VOCÊ CONSUME?}
\end{center}
\par
\textbf{Passo 1:} Quantas horas por dia você fica no celular? \_\_\_\_\_\_ horas
\par
\textbf{Passo 2:} Quanto desse tempo é streaming (Netflix, YouTube, TikTok)? \_\_\_\_\_\_ horas
\par
\textbf{Passo 3:} Multiplique por 30 para saber as horas por mês: \_\_\_\_\_\_ horas/mês
\par
\textbf{Passo 4:} Consumo por hora em diferentes resoluções:
\begin{itemize}
    \item 480p: 0,035 kWh/hora
    \item 1080p: 0,065 kWh/hora
    \item 4K: 0,175 kWh/hora
\end{itemize}
\par
\textbf{Passo 5:} Seu consumo mensal = horas/mês × consumo por hora = \_\_\_\_\_\_ kWh
\par
\textbf{Passo 6:} Custo mensal = consumo × R\$ 0,80 = R\$ \_\_\_\_\_\_
\par
\textbf{Passo 7:} Custo anual = custo mensal × 12 = R\$ \_\_\_\_\_\_
\end{minipage}}%
\lthtmlindisplaymathZ
\lthtmlcheckvsize\clearpage}

{\newpage\clearpage
\lthtmlinlinemathA{tex2html_wrap_inline917}%
\fbox{\begin{minipage}{0.95\textwidth}
\begin{center}\textcolor{dados}{\textbf{\Large CONSUMO DE ENERGIA DOS DATA CENTERS NO BRASIL}}
\end{center}
\par
De acordo com a \textbf{1ª Revisão Quadrimestral das Previsões de Carga para o Planejamento Anual da Operação Energética 2026–2030} (CCEE/ONS/EPE):
\par
\begin{tabular}{|p{4cm}|p{5cm}|}
\hline
\textbf{Indicador} & \textbf{Valor} \\
\hline
Consumo atual de data centers (2026) & 304 MW médios \\
\hline
Consumo previsto (2030) & 3.457 MW médios \\
\hline
\textbf{Crescimento} & \textbf{MAIS DE 11 VEZES em 4 anos} \\
\hline
Investimento previsto até 2030 (Thymos Energia) & R\$ 60 bilhões \\
\hline
Data centers em operação no Brasil & Mais de 130 \\
\hline
\end{tabular}
\par
\vspace{0.5cm}
\par
\textbf{Onde estão os novos data centers?} Regiões Sudeste, Centro-Oeste, Nordeste e Sul.
\par
\textbf{Quantos projetos?} 22 pedidos de acesso à Rede Básica com contratos assinados, sendo 18 já autorizados.
\par
\textbf{Efeito colateral positivo:} Os data centers podem ajudar a diminuir os \textbf{cortes de geração no Nordeste} (usinas eólicas e solares que precisam ser desligadas por excesso de oferta), aumentando o consumo regional.
\par
\vspace{0.3cm}
\textit{Fonte: CCEE/ONS/EPE - 1ª Revisão Quadrimestral 2026-2030; Thymos Energia.}
\end{minipage}}%
\lthtmlindisplaymathZ
\lthtmlcheckvsize\clearpage}

\stepcounter{section}
\stepcounter{section}
\stepcounter{section}
{\newpage\clearpage
\lthtmlinlinemathA{tex2html_wrap_inline943}%
\fbox{\begin{minipage}{0.85\textwidth}
\begin{center}\textbf{\Large O que Christoph Wulf nos ensina}
\end{center}
\par
\begin{tabular}{|p{4cm}|p{5cm}|}
\hline
\textbf{Conceito} & \textbf{O que significa para sua vida} \\
\hline
Antropoceno & A Era dos Humanos – você é responsável pelo futuro do planeta \\
\hline
Nativos digitais & Você nasceu na cultura digital – mas isso não significa que você a entende \\
\hline
Prosumer & Você produz enquanto consome – você é parte da indústria cultural \\
\hline
Teoria Ator-Rede & Máquinas e humanos agem juntos – você não está sozinho(a) na tela \\
\hline
Repetição/Mimesis & Seus hábitos repetitivos te moldam – o que você está repetindo? \\
\hline
Ética da IA & \textit{"Uma nova ética no uso de máquinas é necessária."} \\
\hline
\end{tabular}
\par
\textit{"O que é necessário é uma \textbf{nova ética} no uso de máquinas. As pessoas precisam aprender como a inteligência artificial funciona e como \textbf{integrar princípios éticos} em seu comportamento."} (Wulf, 2022, p. 22)
\end{minipage}}%
\lthtmlindisplaymathZ
\lthtmlcheckvsize\clearpage}

{\newpage\clearpage
\lthtmlinlinemathA{tex2html_wrap_inline947}%
\fbox{\begin{minipage}{0.85\textwidth}
\begin{center}\textcolor{alerta}{\textbf{\Large A PERGUNTA QUE FICA}}
\end{center}
\par
\textit{Se os data centers de IA vão consumir tanta energia quanto o Japão em 2026 e potencialmente o triplo até 2030, quem vai pagar a conta?}
\par
\textit{Os bilhões de usuários que clicam grátis no ChatGPT? As companhias que cobram US\$ 20 por mês? Ou os países que terão de construir usinas inteiras só para sustentar essa nova economia?}
\par
\textit{E o Brasil, que tem energia limpa e sobrando ao meio-dia, vai conseguir capturar essa onda antes que outros países dominem o jogo?}
\end{minipage}}%
\lthtmlindisplaymathZ
\lthtmlcheckvsize\clearpage}


\end{document}
